% Source: https://artofproblemsolving.com/community/c3h1612489_coin_sequence
% Stable identity: aops-topic-1612489-q01
\textbf{Bert chooses THH.} Sam can win only if HHH appears before any T has been flipped, that is, only if the first three flips are HHH. Indeed, once a T has appeared, consider the first time three H's in a row would be completed: the flip just before that run of three H's is a T, so THH has been completed one flip earlier and Bert has already won. Hence
\[
P(\text{Bert wins with THH})=1-\left(\dfrac{1}{2}\right)^3=\dfrac{7}{8} .
\]

\textbf{No choice does better.} Whatever Bert picks, Sam wins whenever the first three flips are HHH, which has probability $\dfrac{1}{8}$; so Bert can never win with probability more than $\dfrac{7}{8}$, and THH achieves this. For comparison, a standard first-step computation gives Bert's winning probabilities $\dfrac{7}{12}$ for THT, $\dfrac{7}{10}$ for TTH, $\dfrac{3}{5}$ for HTT and $\dfrac{1}{2}$ for TTT.

So Bert should choose $\boxed{\textbf{(a)}\ \text{THH}}$, winning with probability $\dfrac{7}{8}$.
